% Guía para presentar el deck OEM (castellano).
% Build:  latexmk -pdf -outdir=build guia-presentacion-oem.tex   (desde paper/oem-brief/)
%
% Acompaña a oem-slides.pdf. Las diapositivas están en inglés: los títulos se
% citan aquí en inglés para localizarlas; la narración va en castellano.
%
% Audiencia: responsable de performance en un OEM de motores, con conocimiento
% de alto nivel y muy crítico con la IA.

\documentclass[11pt,a4paper]{article}

\usepackage[a4paper,top=2.2cm,bottom=2.2cm,left=2.6cm,right=2.6cm]{geometry}
\usepackage[utf8]{inputenc}
\usepackage[T1]{fontenc}
% es-noshorthands: los shorthands de babel-spanish rompen tikz/siunitx.
\usepackage[spanish,es-noshorthands,es-noquoting]{babel}
\usepackage{lmodern}
\usepackage{microtype}
\usepackage{amsmath,amssymb}
\usepackage{booktabs}
\usepackage{xcolor}
\usepackage{enumitem}
\usepackage{fancyhdr}
\usepackage[hidelinks]{hyperref}

\definecolor{inkc}{HTML}{212529}
\definecolor{bluec}{HTML}{4263EB}
\definecolor{redc}{HTML}{A61E4D}
\definecolor{grayc}{HTML}{868E96}
\color{inkc}

\setlength{\parindent}{0pt}
\setlength{\parskip}{0.55em}
\linespread{1.06}
\setlist[itemize]{leftmargin=1.2em,itemsep=2pt,topsep=3pt}

\pagestyle{fancy}
\fancyhf{}
\fancyhead[L]{\footnotesize\color{grayc} Guía de presentación --- deck OEM}
\fancyhead[R]{\footnotesize\color{grayc}\thepage}
\renewcommand{\headrulewidth}{0.2pt}
\renewcommand{\headrule}{\hbox to\headwidth{\color{grayc}\leaders\hrule height \headrulewidth\hfill}}

% \diapo{nº}{título en pantalla}
\newcommand{\diapo}[2]{%
  \vspace{0.9em}%
  \noindent{\color{bluec}\rule{\linewidth}{0.6pt}}\par\vspace{-0.35em}%
  \noindent{\bfseries\color{bluec} Diapositiva #1}\quad{\bfseries #2}\par
  \vspace{-0.15em}}

% narración: lo que se dice en voz alta
\newenvironment{decir}%
  {\begin{list}{}{\leftmargin=1.1em \rightmargin=0em \listparindent=0pt
       \itemindent=0pt \parsep=0.5em \topsep=0.4em}\item[]\ignorespaces}%
  {\end{list}}

% acotación escénica: no se dice
\newcommand{\escena}[1]{{\itshape\color{grayc}[#1]}}
\newcommand{\cifra}[1]{\textbf{#1}}
\newcommand{\idea}[1]{\par\smallskip{\color{grayc}\itshape Idea en una frase: #1}\par}

\title{\vspace{-1.5cm}\textbf{Deck OEM}\\[2mm]
  \large Guía de presentación}
\author{}
\date{\small Acompaña a \texttt{oem-slides.pdf} (17 diapositivas, 2 de reserva)}

\begin{document}
\maketitle
\thispagestyle{fancy}
\vspace{-1.2cm}

\section*{Antes de empezar: la postura}

Tu ventaja es que \textbf{no vienes a vender IA}. Vienes a contar que la mediste y que
casi todo se cayó. Eso desarma al escéptico en el primer minuto, y a partir de ahí te
escucha.

\textbf{Regla de oro:} cada vez que digas algo bueno de la IA, di antes el número que lo
limita. Nunca al revés.

\textbf{Tres cosas que NO debes decir:}
\begin{itemize}
  \item ``La IA revoluciona el mantenimiento'' --- pierdes la sala.
  \item ``Nuestro modelo aprende patrones complejos'' --- suena a folleto.
  \item Cualquier cifra sin decir contra qué se compara.
\end{itemize}

\textbf{Arranque sugerido:}
\begin{decir}
Monté un banco de pruebas para medir si lo que se publica sobre IA en performance se
sostiene. Te adelanto el resultado: la mayoría no. Incluido lo mío. Lo que sí queda es
pequeño y, curiosamente, no es software.
\end{decir}

El texto sangrado es lo que se dice en voz alta; lo que aparece \escena{así} son
acotaciones, no se pronuncian.

\diapo{1}{Portada}
Léela y pasa. No te detengas.

\diapo{2}{Start from the skeptical position}
\begin{decir}
La pregunta era simple: ¿la IA gana al método clásico de análisis de gas path, o no? La
trampa habitual es comparar una IA bien ajustada contra un método clásico mal ajustado.
Aquí les di a los dos el mismo dato, la misma métrica y el mismo presupuesto de ajuste. Y
escribí las hipótesis \textbf{antes} de correr nada, para no poder cambiarlas después.

\escena{señala la columna derecha} Ahí tienes lo que salió. Un paper publicado cuyo
titular resulta ser ruido. Dos titulares míos que se caen a la tercera parte. Y la tarea
donde más se vende la IA acaba en empate.
\end{decir}
\idea{no vengo a convencerte, vengo a enseñarte una auditoría.}

\diapo{3}{The rig, in one slide}
\begin{decir}
Construí un motor virtual del CFM56-7B calibrado contra datos de certificación públicos,
y con él una flota de cien motores que viven hasta que salen de ala. La clave es esta:
\textbf{de cada motor sé la verdad}. Sé exactamente cuánto se ha degradado cada
componente en cada vuelo.

Eso con datos reales no lo puedes hacer. Con datos reales solo puedes comparar una
estimación contra otra estimación. Aquí puedo preguntar lo que de verdad importa: ¿el
diagnóstico \textbf{acertó}?
\end{decir}

\escena{si pregunta ``¿y eso es realista?''}
\begin{decir}
El modelo está calibrado contra la hoja de certificación del motor y contra el banco de
datos de emisiones de OACI, y predice el consumo de combustible dentro de un \cifra{3\,\%}
de las medidas de banco certificadas. Los números absolutos no los defiendo. Lo que se
transfiere es la \textbf{geometría} --- cuántas incógnitas hay contra cuántas medidas, y
cómo de parecidas son las huellas de dos averías --- porque eso depende de cómo está
instrumentado el motor, no de si mi modelo acierta el consumo al medio por ciento.
\end{decir}

\diapo{4}{A peer-reviewed deep-learning result, replicated exactly}
\escena{esta es la diapositiva importante; tómate tiempo}
\begin{decir}
Cogí un paper de 2026, publicado en Springer, revisado por pares. Dice que una red
neuronal concreta es la mejor de quince arquitecturas para predecir vida remanente. Lo
reimplementé exactamente como está escrito.

Primero, lo justo: \textbf{replica}. La métrica es el error medio al predecir cuántos
ciclos le quedan al motor antes de salir de ala. Ellos publican \cifra{14,12 ciclos} de
error; a mí me sale \cifra{14,56 ciclos}. No estoy montando un hombre de paja.

Ahora el problema. Ellos ordenan quince arquitecturas con diferencias de entre
\cifra{0,12 y 0,60 ciclos} entre una y la siguiente. Pero si entrenas \textbf{el mismo
modelo dos veces}, cambiando solo el número aleatorio con el que arranca el entrenamiento,
el resultado se mueve \cifra{2,94 ciclos}. Es decir: la distancia entre modelos distintos
es cinco veces más pequeña que la dispersión de repetir el mismo modelo. Ese ranking no
mide la arquitectura, mide el azar. Y ellos publican una sola corrida por modelo.

Y hay algo peor \escena{señala los puntos rojos}: en \cifra{cinco de diez} corridas el
modelo no llega a entrenar. Se queda prediciendo cero. Y la causa está en sus propias
tablas: una configuración que, si el modelo se atasca al principio, no puede recuperarse.
\end{decir}

\escena{analogía, si hace falta}
\begin{decir}
Es como ordenar quince motores en banco por consumo específico, con una décima de
diferencia entre uno y el siguiente, cuando repetir el ensayo del mismo motor te da tres
décimas de dispersión. El orden que publicas es el del ruido del banco.
\end{decir}

\diapo{5}{The evaluation question to ask anyone}
\begin{decir}
De todo esto, si te llevas una sola cosa, que sea esta pregunta: \textbf{¿cuál es la
dispersión entre repeticiones?} Antes de mirar el titular.

Sirve para un proveedor que te enseña un piloto. Sirve para un equipo interno. Y sirve
para mí --- de hecho me la apliqué, y es la siguiente diapositiva.
\end{decir}
\idea{un ranking dentro del ruido no es un resultado.}

\diapo{6}{We applied it to our own claims, and both shrank}
\begin{decir}
Me apliqué el mismo criterio y perdí dos titulares.

El primero: yo decía que la IA daba intervalos de incertidumbre \cifra{seis veces} más
estrechos. Falso tal como estaba dicho. Mi intervalo estaba \emph{calibrado} --- ajustado
para que cuando digo noventa por ciento, sea de verdad noventa. El del método clásico no
lo estaba, y por eso salía anchísimo. Calibras los dos igual y la ventaja se queda en
\cifra{1,34}. Tres cuartas partes de aquel seis era la calibración, no el modelo. Y la
calibración es una técnica que el lado clásico puede coger gratis.

El segundo: ``la IA predice la vida remanente entre 2,3 y 4,4 veces mejor''. Cierto, pero
solo contra la extrapolación lineal que se usa hoy en operación. Contra un método clásico
bueno es \cifra{1,34} veces al noventa por ciento de la vida, y \textbf{empate} al
setenta.
\end{decir}
\idea{cuando la línea base es competente, la ventaja es del 30\,\%, no de varias veces.}

\diapo{7}{Where AI lost outright}
\begin{decir}
La tarea es: ha aparecido una degradación, y hay que decir \textbf{en qué componente}
está. Sobre los casos difíciles --- los pares de componentes que dejan una señal parecida
--- el método clásico acierta el \cifra{31\,\%} de las veces. La IA acierta el
\cifra{31\,\%}. Idéntico.

Y esto no es que el algoritmo sea malo. Es que \textbf{la información no está en el
dato}. En un informe de cabina tienes tres medidas que sirven --- \textbf{N2, caudal de
combustible y EGT} --- y diez parámetros de salud que estimar, rendimiento y capacidad de
cada módulo. Es como pedirme tres números que sumen diez: hay infinitas combinaciones y
ninguna manera de elegir entre ellas.

\escena{esto es el corazón del argumento; explícalo despacio} Y hay algo más fino. Cada
avería deja una huella con forma propia: cuánto sube la EGT, cuánto sube el combustible,
cuánto cambia N2. Si dos averías mueven los tres sensores en proporciones distintas, las
distingues. Si los mueven casi igual, no. Esa diferencia de forma se puede medir como un
ángulo: \cifra{90 grados} sería perfectamente distinguibles, \cifra{0 grados} sería
idénticas. Entre el rendimiento del compresor de alta y el de la turbina de alta hay
\cifra{1,3 grados}. Prácticamente la misma huella.

Lo comprobé por cuatro caminos independientes. Reduje a la mitad el ruido de todos los
sensores: el acierto se queda en \cifra{31--38\,\%}, no se mueve. Metí un ``sensor
virtual'' calculado a partir de los que ya hay, que es lo que venden algunos proveedores:
\cifra{cero mejora}, exactamente cero, porque no añade información nueva. Barrí los seis
puntos de operación calibrados --- despegue, subida, tres cruceros --- por si el ángulo se
abría en alguno: el mejor da \cifra{1,55 grados} y juntándolos todos sale
\cifra{1,38}. Y ningún método aprendido lo movió.
\end{decir}

\escena{remate, pausa antes}
\begin{decir}
Ningún algoritmo arregla que falte información.
\end{decir}

\diapo{8}{And most prognostic error is not a method gap}
\begin{decir}
Y una última mala noticia, que en realidad es una decisión de presupuesto. Cogí motores
que están en el \emph{mismo estado real de degradación} --- eso lo sé porque en mi flota
conozco la verdad --- y miré cuánto difiere su vida restante de verdad. Resultado: a mitad
de vida, el \cifra{87\,\%} de la dispersión en vida remanente es \textbf{irreducible}.
Dos motores idénticos en condición fallan con cientos de ciclos de diferencia por cosas que
ninguna medida de hoy contiene.

Traducido: a mitad de vida, el error del mejor método está a \cifra{1,6 veces} de ese
suelo físico. Queda muy poco que ganar. Al final de la vida está a \cifra{4 veces} del
suelo: ahí sí hay margen real, y ahí sí paga financiar un algoritmo mejor.
\end{decir}

\diapo{9--10}{An instrumentation design tool / Why this one is ours to act on}
\escena{cambio de tono: llevas media presentación diciendo que no; ahora el sí}
\begin{decir}
Lo más sólido que salió de todo esto \textbf{no es IA}. Es teoría de estimación clásica
aplicada a través del modelo del motor.

La idea: para un motor y sus condiciones reales de vuelo, puedo calcular cuánta
información sobre cada componente llega de verdad a los sensores. Y de ahí sale un
límite: esto lo puedes saber, esto no lo puede saber \emph{nadie} con esos sensores.

\escena{si te preguntan si esto está validado, di el matiz antes de que lo busquen} Una
cosa que encontramos tarde y nosotros mismos: la afirmación de que esa cota \emph{ordena} el
error real del estimador ---correlación 0,70--- no supera su propio control. Un cuarto de ese
0,70 es la matriz que la cota comparte con el estimador; la mitad sobrevive si usas un
estimador que no ve esa matriz; y ninguna de las dos llega a significación, porque con diez
direcciones un test de rangos no puede. No está demostrado que la cota falle: está demostrado
que no se puede probar como la probamos.

Y si te aprietan más, tienes la respuesta buena: sabemos exactamente qué haría falta para
probarla. La cota es \emph{por motor}, así que hace falta una flota donde los motores vuelen
misiones de verdad distintas. En la nuestra esa magnitud varía un 1,15\,\%: todos vuelan
casi igual. Por eso la frase precisa no es «sin probar» sino \textbf{«no medible en esta
flota»}, que es una posición mucho más fuerte ---dice qué experimento la cierra--- y además
es honesta. Las cifras de \emph{adquisición} que vienen ahora son medidas distintas y no
están afectadas.

Mira la tabla. Si añades sondas de estación --- presión y temperatura a la salida del
fan, a la salida del compresor de alta --- la incertidumbre sobre el rendimiento del
compresor de alta se reduce \cifra{45 veces}: pasa de un valor que no te deja concluir
nada a uno con el que decides. Y el acierto en aquellos casos difíciles, el mismo 31\,\%
de antes, pasa de \cifra{15\,\% al 92\,\%}. Si en cambio reduces el ruido de las
sondas que ya tienes: no cambia nada. Si compras software con ``sensores virtuales'': no
cambia nada.

Y aquí está la consecuencia para nosotros: cualquier servicio de EHM montado \textbf{solo
sobre datos de cabina} tiene un techo que ningún roadmap de software levanta. La salida es
instrumentación. Y la instrumentación la decidimos nosotros, no el cliente.
\end{decir}
\idea{esto pone precio a una decisión de hardware antes de comprometer el hardware.}

\diapo{11}{The one AI result that held up}
\begin{decir}
Y ahora sí, lo único de IA que aguantó, y es más estrecho y más útil que lo que se suele
contar.

Meter física dentro del modelo --- darle al aprendiz la estimación del estado del motor
--- \textbf{empeora} la predicción de vida remanente. Lo comprobé dos veces y las dos
falló. Pero esa misma inyección \textbf{mejora claramente} decir qué mecanismo se está
comiendo el motor: la varianza explicada sube \cifra{0,14} sobre un máximo de 1, con una
probabilidad de que sea casualidad de \cifra{4 entre diez mil}.

Y lo monté con un brazo de control, para que no me pudiera engañar a mí mismo: si el
control también hubiera ganado, entonces la mejora sería solo ``más datos de entrada''.

La explicación es mecánica y tiene sentido: la vida remanente ya va en el margen de EGT
que el modelo está viendo, así que el estado físico no le añade nada. Pero \emph{qué
módulo} se está degradando es justo lo que ese estado codifica.

Conclusión: ``physics-informed'' no es ni magia ni humo. \textbf{Depende de la tarea.} Y
ahora sé de cuál.
\end{decir}

\diapo{12}{What this is, and what it is not}
\escena{dila entera, sin acelerar: es la que te da credibilidad}
\begin{decir}
Qué \textbf{no} es esto. No es un producto. No es una herramienta validada para uso
operativo. Y los números absolutos no los defiendo: la flota es sintética y el modelo está
calibrado contra datos públicos, no contra nuestro banco.

Qué \textbf{sí} es: un método reutilizable para decidir qué afirmaciones sobre IA en
performance financiamos y cuáles rechazamos. Con la disciplina que hace que las respuestas
valgan algo --- hipótesis congeladas antes de correr, mismo presupuesto de ajuste para los
dos lados, los resultados negativos con el mismo peso que los positivos, y ninguna cifra
escrita a mano: todas salen del código.
\end{decir}

\diapo{13}{Proposed next step}
\begin{decir}
La propuesta es una sola: repetir el análisis de instrumentación sobre \textbf{nuestros}
modelos de motor y sobre las configuraciones de sensores que estemos considerando.

No necesita IA. No necesita recoger datos nuevos. Y responde a una pregunta que ya estamos
pagando por responder de otra manera.
\end{decir}
\escena{y cállate: deja que hable él}

\diapo{14}{Reserva --- The discipline that makes the answers usable}
\escena{sácala si pregunta ``¿y cómo sé que no te has engañado a ti mismo?''; es la
pregunta correcta y conviene que la haga}
\begin{decir}
Cinco cosas, y cada una tiene una cicatriz detrás.

\textbf{Una: escribir la hipótesis antes de correr.} Antes de cada experimento dejé un
documento fechado con qué iba a medir y qué umbral contaba como éxito. Cuando el resultado
salió en contra, quedó en contra. La hipótesis de que la IA aislaría mejor las averías está
refutada en el informe, con su número, y no la he tocado.

\textbf{Dos: el mismo presupuesto de ajuste para los dos lados}, y aquí está la cicatriz
más útil. En mi primera versión el detector de IA sacaba un acierto de \cifra{0,06} ---
inservible --- y el clásico le ganaba de calle. Con el mismo número de intentos de ajuste
para ambos, la IA subió a \cifra{0,48} y el marcador \textbf{se dio la vuelta}. O sea: si
me hubiera quedado en la primera comparación habría publicado exactamente la conclusión
contraria. Por eso, cuando alguien te enseñe un marcador, la pregunta es si las dos partes
estaban igual de ajustadas.

\textbf{Tres: puertas que pueden fallar, y fallaron.} Puse una condición a mi propia flota
sintética: un clasificador tonto, de los que no aprenden nada, no debe pasar del 60\,\% de
acierto. Si lo pasa, es que el problema es demasiado fácil y cualquier método parecerá
bueno. La primera flota dio \cifra{81,6\,\%}. \textbf{Regeneré la flota, no bajé el
umbral.} Acabó en \cifra{54\,\%}, y falló dos veces por el camino.

\textbf{Cuatro: trampas para cazarme a mí mismo.} Ejemplo real: metí en la flota ocho
motores con una deriva en un sensor que los datos de cabina \emph{físicamente no pueden
ver}. Si mi método hubiera dicho que los detecta, sería que está leyendo otra cosa, y todo
el resultado sería basura. No los detectó, como debía. Eso es lo que me permite afirmar
que, cuando digo que algo no se puede distinguir, es que de verdad no se puede.

\textbf{Cinco: ninguna cifra escrita a mano.} Todas las gráficas y tablas se regeneran
desde el código. Si un número cambia, cambia en el documento solo.
\end{decir}
\idea{la credibilidad no está en el resultado, está en haber podido perder.}

\diapo{15}{Reserva --- The wall is a snapshot property}
\escena{sácala si pregunta ``¿entonces con los sensores de cabina no se puede diagnosticar
nada?''}
\begin{decir}
No, hay un matiz importante, y es el resultado más fino que salió.

El muro del que hablé es sobre \textbf{una foto}. En un informe suelto no puedes separar
si el rendimiento que se ha perdido es del compresor de alta o de la turbina de alta.

Pero un motor no es una foto, es una película de miles de vuelos. Y ahí sí aparece
información, porque los mecanismos de degradación tienen \textbf{formas distintas en el
tiempo}. El ensuciamiento del compresor se recupera en parte con cada lavado: hace dientes
de sierra. El asentamiento de holguras ocurre rápido al principio y luego para: hace un
codo. El deterioro de la sección caliente se acelera al final de la vida. Son firmas
temporales que no se parecen en nada, aunque en una foto suelta sí se parezcan.

Midiéndolo: de cinco mecanismos, \textbf{tres} se pueden atribuir a partir de la forma de
la trayectoria. Y el mejor recuperado, con \cifra{0,78} de varianza explicada sobre 1, es
justamente el de la sección caliente, que es el que provoca la confusión en la foto.

La regla que salió de todo esto: \textbf{la forma de la trayectoria separa lo que tiene
forma}. Un diente de sierra y un codo se leen. Una rampa lisa, no --- por mucho que
esperes. Y por eso mismo tampoco se puede distinguir un sensor que va derivando poco a
poco de una degradación real: las dos son rampas lisas.
\end{decir}

\escena{si pregunta ``¿y eso para qué me sirve operativamente?''; contesta con honestidad,
porque el límite es real}
\begin{decir}
Para decidir \emph{cuánto} de la pérdida es lavable y cuánto es permanente, sí sirve: eso
se estima bien. Para decidir \emph{qué módulo} tocar en el taller, en esta flota no sirve,
porque la sección caliente domina en los cien motores de cien. Es una limitación de mi
flota sintética, no un resultado sobre motores reales, y está declarada como tal en el
informe.
\end{decir}
\idea{el muro es de la foto, no de la película --- pero solo si en la película pasa algo.}

\section*{Objeciones probables, y qué contestar}

\textbf{``Esto es todo simulado, no me sirve.''}
\begin{decir}
Correcto, y por eso no defiendo ningún número absoluto. Lo que se transfiere es la
geometría: cuántas incógnitas hay contra cuántas medidas. Eso es igual en nuestro motor
real, y es lo que manda. Además, el análisis de instrumentación corre sobre nuestros
modelos, donde es más fuerte que aquí.
\end{decir}

\textbf{``La IA es humo.''}
\begin{decir}
En tres de las cuatro tareas que medí, te doy la razón con números: aislamiento de
averías es empate exacto, la incertidumbre calibrada baja de seis veces a
\cifra{1,34}, y la física metida en el modelo empeoró la predicción de vida dos veces
seguidas. En la cuarta, predicción de vida, hay un \cifra{34\,\%} de mejora contra un
método clásico bueno, y solo en el último tramo de vida. Es real y es pequeño. Y lo que te
propongo no depende de que la IA funcione.
\end{decir}

\textbf{``¿Cuánto ha costado esto?''}
\begin{decir}
El estudio principal completo --- generar la flota, correr los dos métodos, sacar los
veredictos --- son \cifra{25 minutos} en un portátil. Los estudios de extensión más
pesados, entre media hora y hora y media cada uno usando la GPU del portátil. Sin
infraestructura. Y todo viene de datos públicos de certificación, así que no hay problema
de propiedad intelectual para circularlo internamente.
\end{decir}

\textbf{``¿Por qué no lo ha hecho ya un proveedor?''}
\begin{decir}
Porque hace falta conocer la verdad del motor para validarlo, y eso solo lo tienes en
simulación. Un estudio con datos reales no puede comprobar si su propia estimación de
incertidumbre es honesta. Nosotros sí.
\end{decir}

\textbf{``¿Y si el proveedor X dice que su modelo da un 95\,\%?''}
\begin{decir}
Primera pregunta: ¿contra qué línea base, y estaba ajustada igual? Segunda: ¿cuál es la
dispersión entre repeticiones? Con esas dos preguntas se cae la mayoría.
\end{decir}

\section*{Cifras clave, con unidades}

Si te preguntan por una cifra y dudas, está aquí. Ninguna va sin su unidad y sin contra
qué se compara.

\begin{center}\small
\begin{tabular}{@{}p{0.40\textwidth}p{0.26\textwidth}p{0.28\textwidth}@{}}
  \toprule
  Cifra & Unidad / qué mide & Contra qué \\
  \midrule
  $14{,}56$ frente a $14{,}12$ & ciclos de error medio al predecir vida remanente & el paper publicado, misma métrica \\
  $0{,}12$--$0{,}60$ & ciclos entre una arquitectura y la siguiente & --- \\
  $2{,}94$ & ciclos de dispersión al repetir \emph{el mismo} modelo & el ruido que invalida el ranking \\
  $5$ de $10$ & corridas que no llegan a entrenar & --- \\
  31\,\% vs 31\,\% & acierto identificando el componente degradado, casos difíciles & clásico vs IA \\
  $1{,}3^{\circ}$ & separación entre dos huellas de avería, sobre $90^{\circ}$ posibles & $90^{\circ}$ sería distinguirlas perfectamente \\
  $1{,}38^{\circ}$ & lo mismo, usando los seis puntos de operación juntos & frente a $1{,}55^{\circ}$ solo en crucero \\
  $45\times$ & reducción de la incertidumbre sobre el rendimiento del compresor de alta & sondas de estación vs solo cabina \\
  15\,\% $\to$ 92\,\% & el mismo acierto de antes, con sondas de estación & vs el 31\,\% de cabina \\
  87\,\% & parte de la dispersión en vida remanente que es irreducible, a mitad de vida & medido contra la verdad de la flota \\
  $1{,}6\times$ / $4\times$ & distancia del mejor método al suelo físico, a media y a final de vida & el suelo es lo irreducible \\
  $6\times \to 1{,}34\times$ & estrechez del intervalo de incertidumbre & antes y después de calibrar los dos lados igual \\
  $1{,}34\times$ & error en vida remanente & IA vs método clásico bueno, al 90\,\% de vida \\
  empate & lo mismo, al 70\,\% de vida & --- \\
  $+0{,}14$ & varianza explicada al atribuir mecanismo, sobre un máximo de $1$ & con física inyectada vs sin ella \\
  $25$ min & estudio principal completo & en un portátil, sin infraestructura \\
  \bottomrule
\end{tabular}
\end{center}

\textbf{Si no te sabes una cifra, no la inventes.} Di ``eso lo tengo medido, te lo paso'',
y sigue. Con este interlocutor, una cifra dudosa cuesta más que no darla.

\section*{Si solo tienes cinco minutos}

Diapositivas \textbf{4, 6, 7, 9 y 13}, en ese orden.

Un paper publicado que no se sostiene $\rightarrow$ mis propios números corregidos
$\rightarrow$ dónde la IA pierde y por qué $\rightarrow$ lo que sí sirve, que es
instrumentación $\rightarrow$ la propuesta.

\end{document}
